\section{series with variable sign}
\label{sec:serie_segno_variabile}
\index{series!alternating signs}

The series $\sum \frac{(-1)^k}{k+1}$ whose sum can be expressed as
\mymargin{alternating harmonic series}%
\index{series!harmonic!alternating}%
\[
1 - \frac{1}{2} + \frac{1}{3} - \frac{1}{4} +  \frac{1}{5} \dots
\]
is not absolutely convergent
(since the series $\sum \frac 1 {k+1}$ is divergent) but has a generic
infinitesimal term.
Therefore, we do not have any criterion that allows us to determine its nature.
However, we can exploit the fact that the signs are \emph{alternating}, meaning
that the terms with even indices have opposite signs to those with odd indices.
Indeed, it is noted that setting
\[
  S_n = \sum_{k=0}^n \frac{(-1)^k}{k+1}
\]
we have
\begin{align*}
S_{2n+2}
  &= S_{2n+1} + \frac{1}{2n+3}
  = S_{2n} - \frac{1}{2n+2} + \frac{1}{2n+3}
  < S_{2n}\\
S_{2n+3}
  &= S_{2n+2} - \frac{1}{2n+4}
  = S_{2n+1} + \frac{1}{2n+3} - \frac{1}{2n+4}
  > S_{2n+1} \\
\end{align*}
Thus, the sequence of partial sums with even indices is decreasing while
the sequence with odd indices is increasing. Therefore, both
subsequences have a limit: $S_{2n} \to S$, $S_{2n+1} \to R$.

But
\[
    S - R = \lim_{n\to +\infty} (S_{2n} - S_{2n+1}) =
  \lim_{n\to+\infty}\frac{1}{2n+2} = 0.
\]
Thus $S=R$ and the entire sequence has limit $S$.
On the other hand, $S \le S_0$ since $S_{2n}$ is decreasing and $S\ge S_1$ since
$S_{2n+1}$ is increasing.
We conclude that $S$ is finite and thus the series is convergent.%
\mynote{%
To determine the value of the sum of this series, more advanced tools are
needed.
See equation~\ref{eq:serie_ln2}.
}

This proof can be generalized in the following.

\begin{theorem}[Leibniz criterion for alternating series]
\label{th:Leibniz}%
\mymark{***}%
\mymargin{Leibniz criterion}%
\index{criterion!Leibniz}%
\index{theorem!Leibniz}%
\index{Leibniz!criterion}%
Let $b_n$ be a monotonic and infinitesimal sequence. Then
the series
\[
  \sum (-1)^{n} b_n
\]
is convergent.

More precisely, if $\displaystyle S_n = \sum_{k=0}^n (-1)^k b_k$
are the partial sums,
it is observed that the sum of the series $S= \lim S_n$ is always between
two consecutive terms of the sequence $S_n$.
That is, if $b_k\ge 0$, we have
\[
  S_{2n+1} \le S \le S_{2n}.
\]
\end{theorem}
%
\begin{proof}
\mymark{***}
Without loss of generality, we can assume that the sequence $b_n$ is decreasing
and thus $b_n \ge 0$ (since the limit is zero).
Setting
\[
 S_n = \sum_{k=0}^n (-1)^k b_k
\]
we have
\begin{align*}
  S_{2n+2} &= S_{2n} - b_{2n+1} + b_{2n+2} \\
  S_{2n+3} &= S_{2n+1} + b_{2n+2} - b_{2n+3}.
\end{align*}
Since $b_n$ is decreasing, we have $b_{2n+2} < b_{2n+1}$ and $b_{2n+3} <
b_{2n+2}$ from which
\[
  S_{2n+2} < S_{2n}, \qquad S_{2n+3} > S_{2n+1}.
\]
Thus, the sequences $S_{2n}$ and $S_{2n+1}$ are monotonic and consequently
have a limit:
\[
  S_{2n} \to S, \qquad S_{2n+1} \to R
\]
with $S, R  \in [-\infty, +\infty]$.
Moreover, since $b_n$ is infinitesimal
\[
  S - R
  = \lim_{n\to +\infty} S_{2n} - S_{2n+1}
  = \lim_{n\to +\infty} b_{2n+1} = 0.
\]
Thus $S=R$. Furthermore, since $S_{2n}$ is decreasing, we have
$S \le S_0$ and since $S_{2n+1}$ is increasing, we have $S\ge S_1$.
Thus $S$ is finite and the series converges.

We have also obtained that
$S_{2n-1} \le S \le S_{2n}$ and $S_{2n+1} \le S \le S_{2n}$
thus the second part of the statement is also verified.
\end{proof}

The previous theorem applies, for example, to the following series,
which are convergent but not absolutely convergent:
\[
   \sum \frac{(-1)^n}{\sqrt n}, 
   \sum \frac{(-1)^n}{n\ln n}.
\]

Care should be taken to verify the hypotheses in the previous theorem.
Since it is an alternating series, the asymptotic comparison criterion cannot be
applied to reduce the terms $b_n$ to the terms of a simpler series to study.
In particular, the monotonicity of the terms $b_n$ is not guaranteed by the
existence of an asymptotic and monotonic sequence.
See the following.

\begin{example}
Let $a_k = \frac{(-1)^k}{\sqrt k} - \frac{1}{k}$.
Note that $a_k \sim c_k = \frac{(-1)^k}{\sqrt k}$
but the two series $\sum a_k$ and $\sum c_k$ have different natures:
the first is divergent as it breaks into the difference 
of a convergent series and a divergent one,
while the second is convergent by the Leibniz criterion
(theorem~\ref{th:Leibniz}).
\end{example}

Thanks to the Leibniz criterion, it is easy 
to produce examples of convergent series
but not absolutely convergent.
The following theorem tells us that for series of this type 
the sum of the series depends on the order in which the terms are taken:
indeed, any desired value can be obtained as the sum.

\begin{theorem}[conditional convergence]%
\label{th:convergenza_condizionata}%
\index{convergence!conditional of a series}%
\mymargin{conditional convergence}%
Let $\sum a_k$ be a convergent but not absolutely convergent series with real
terms.
Then, given any $x \in [-\infty , +\infty]$, there exists a bijective
rearrangement
$\sigma \colon \NN \to \NN$ such that
\[
  \sum_{k=0}^{+\infty}  a_{\sigma(k)} = x.
\]
\end{theorem}
%
\begin{proof}
We divide the terms of the sequence $a_k$ into non-negative terms
and negative terms. Let $a^+_k$ be the subsequence of non-negative terms
and $-a^-_k$ the subsequence
of negative terms (thus $a^+_k\ge 0$ and $a^-_k > 0$). We will have
\begin{align*}
  \sum_{k=0}^n a_k &= \sum_{k=0}^{n^+} a^+_k - \sum_{k=0}^{n^-} a^-_k \\
  \sum_{k=0}^n \abs{a_k} &= \sum_{k=0}^{n^+} a^+_k + \sum_{k=0}^{n^-} a^-_k
\end{align*}
where $n^+ +1$ and $n^-+1$ are respectively
the number of non-negative and negative terms
among the first $n+1$ terms of the sequence $a_k$.

Now observe that it must be
\[
\sum_{k=0}^{+\infty} a_k^+ = +\infty \qquad \text{and} \qquad
\sum_{k=0}^{+\infty} a_k^- = +\infty.
\]
First, the sums exist because the series are non-negative.
If both these sums were finite, then the series $\sum\abs{a_k}$ would be
convergent, but by hypothesis, we have assumed that $\sum a_k$ is not
absolutely convergent.
Therefore, at least one of the two sums is infinite. If the sum of the positive
terms
were infinite and that of the negative terms were finite, we could conclude that
the sum of the series $\sum a_k$ would be infinite.
Conversely, if the sum of the positive terms were finite and that of the
negative terms were infinite, the sum $\sum a_k$ would be $-\infty$. But by
hypothesis,
we have required that the series $\sum a_k$ be convergent.

Given $x\in \RR$, we can then start summing the positive terms
$a^+_k$ until we reach or exceed the value $x$. At that point, we begin summing
the negative terms until we fall below the value $x$.
Then we continue summing the positive terms until we exceed $x$
and again continue with the negative terms until we fall below $x$. By
appropriately interspersing positive and negative terms,
we can thus obtain partial sums that oscillate around the value 
of $x$ and get closer and closer to $x$ since at each change of "direction" 
the distance from $x$ is less than the absolute value of the last term summed 
and the sequence of terms $a_k$ is infinitesimal since the series $\sum a_k$ is
convergent.

The same can be done to obtain a sum $x=+\infty$. Given any
sequence $x_n \to +\infty$, start summing the positive terms until you exceed
the value $x_1+a_1^-$.
Then sum a single negative term, $-a_1^-$, and obtain a sum greater than $x_1$. 
Then sum enough positives until you exceed $x_2+a_2^-$. 
Then sum another single negative term and so on. 
Clearly, the sums will tend to $+\infty$.

The case $x=-\infty$ is treated similarly.
\end{proof}